\documentclass[10pt, aspectratio=169]{beamer}

\usepackage[spanish]{babel}
\usepackage[utf8]{inputenc}
\usepackage[pscoord]{eso-pic}
\usepackage{hyperref}
\usepackage{eqnarray,amsmath}
\usepackage[absolute,overlay]{textpos}
\hypersetup{colorlinks=true,urlcolor=blue}

\usepackage{tikz}
\usetikzlibrary{automata, positioning, arrows.meta, petri}

\usepackage{listings}
\usepackage{xcolor}

\usepackage[style=apa]{biblatex}
\bibliography{biblio.bib}

\title{}
\author{}
\date{}

\newtheorem{teorema}{Teorema.}
\setbeamersize{text margin left=10pt,text margin right=10pt}
\beamertemplatenavigationsymbolsempty

\hypersetup{
	colorlinks=true,
	pdfborderstyle={/S/U/W 1}
}

\setbeamertemplate{background}{\includegraphics[width=\paperwidth,height=\paperheight,keepaspectratio]{portada.pdf}}

\newcommand{\azulULA}{\color[rgb]{0.30078125,0.30078125,0.999999}}
\definecolor{azulULA}{RGB}{77, 77, 255}

\setbeamercolor{title}{fg=azulULA}
\setbeamercolor{frametitle}{fg=azulULA}
\setbeamercolor{block title}{bg=azulULA,fg=white}
\setbeamercolor{block body}{bg=azulULA!10!bg}

\renewcommand{\qedsymbol}{\textcolor{black}{\openbox}}

\newcommand{\placetextbox}[3]{
	\setbox0=\hbox{#3}
	\AddToShipoutPictureFG*{
		\put(\LenToUnit{#1\paperwidth},\LenToUnit{#2\paperheight}){\vtop{{\null}\makebox[0pt][c]{#3}}}
	}
}
\newcommand{\placetextboxleft}[3]{
	\setbox0=\hbox{#3}
	\AddToShipoutPictureFG*{
		\put(\LenToUnit{#1\paperwidth},\LenToUnit{#2\paperheight}){\vtop{{\null}\makebox[0pt][l]{#3}}}
	}
}

\begin{document}
	% Diapositiva de portada
	\begin{frame}
		\placetextbox{0.47}{0.40}{\large\azulULA\textbf{\begin{tabular}{c} Teoría de la Computación y Lenguajes Formales: \\
					Sesión 12-A. Verificación Formal de Workflows: El Análisis Previo \end{tabular}}}
	\end{frame}
	
	% Cambio de formato para el cuerpo
	\setbeamertemplate{background}{
		\includegraphics[width=\paperwidth,height=\paperheight,keepaspectratio]{fondo1280x720.pdf}}
	
	\begin{frame}{Alcance de la sesión 12-A}
		\tableofcontents
	\end{frame}

% =======================================================
\section{La Necesidad de la Verificación Previa}
% =======================================================

\begin{frame}{}
	\centering
	\Huge{La Necesidad de la Verificación Previa: \\ ¿Por qué no probar en producción?}
\end{frame}

\begin{frame}{La Falacia de las Pruebas en Sistemas Concurrentes}
	\begin{block}{El problema del No-Determinismo y la Concurrencia}
		En sistemas secuenciales clásicos, probar un conjunto representativo de casos de prueba puede ser suficiente. En \textbf{sistemas concurrentes (fábricas, flujos de trabajo administrativos)}:
		\begin{itemize}
			\item Existen millones de posibles intercalaciones de eventos (\emph{interleavings}).
			\item Los errores de concurrencia (\emph{deadlocks}, condiciones de carrera) son \textbf{intermitentes} y casi imposibles de reproducir mediante pruebas empíricas.
			\item \alert{Probar en producción es inadmisible:} parar una celda robótica, duplicar un pago bancario o dejar pedidos atrapados tiene un costo inmenso.
		\end{itemize}
	\end{block}

	\begin{alertblock}{La Promesa de las Redes de Petri}
		Una especificación formal en Redes de Petri permite \textbf{demostrar matemáticamente} que el sistema jamás se bloqueará, que no desbordará recursos y que todo caso culminará correctamente, \textbf{antes de conectar un sensor o compilar una sola línea de código}.
	\end{alertblock}
\end{frame}

\begin{frame}{El Ciclo de Vida Completo de la Ingeniería de Procesos}
	\begin{block}{Metodología de 8 Pasos: De los Requisitos a la Operación}
		Cada tema de la materia representa una compuerta estricta de validación formal:
	\end{block}

	\begin{center}
		\begin{tikzpicture}[x=3.4cm, y=-1.8cm, >=Stealth, thick]
			% Fila 1: Especificación y Sintaxis
			\node[draw=azulULA, rectangle, rounded corners, fill=azulULA!10, text width=2.6cm, align=center] (p1) at (0, 0)
				{\textbf{1. Entender}\\ \scriptsize Requisitos del dominio\\ y restricciones};
				
			\node[draw=azulULA, rectangle, rounded corners, fill=azulULA!10, text width=2.6cm, align=center] (p2) at (1, 0)
				{\textbf{2. Modelar RdP}\\ \scriptsize Lugares, transiciones\\ y recursos};
				
			\node[draw=azulULA, rectangle, rounded corners, fill=azulULA!10, text width=2.6cm, align=center] (p3) at (2, 0)
				{\textbf{3. Codificar}\\ \scriptsize Formalismo\\ RdP + Ontología};
				
			\node[draw=orange!80!black, rectangle, rounded corners, fill=orange!10, text width=2.6cm, align=center] (p4) at (3, 0)
				{\textbf{4. Validar (PDA)}\\ \scriptsize Sintaxis / Gramática\\ (\textbf{Clase 10})};

			% Fila 2: Semántica, Verificación y Run-Time
			\node[draw=orange!80!black, rectangle, rounded corners, fill=orange!10, text width=2.6cm, align=center] (p5) at (3, 1)
				{\textbf{5. Completitud}\\ \scriptsize Semántica Ontología\\ y recursos tipados};
				
			\node[draw=red!80, rectangle, rounded corners, fill=red!15, text width=2.6cm, align=center] (p6) at (2, 1)
				{\textbf{6. Verificar}\\ \scriptsize Árbol de alcance,\\ \emph{Soundness} (\textbf{12-A})};
				
			\node[draw=green!60!black, rectangle, rounded corners, fill=green!15, text width=2.6cm, align=center] (p7) at (1, 1)
				{\textbf{7. Desplegar}\\ \scriptsize Inyección al Motor\\ Fénix};
				
			\node[draw=azulULA, rectangle, rounded corners, fill=azulULA!25, text width=2.6cm, align=center] (p8) at (0, 1)
				{\textbf{8. Usar Sistema}\\ \scriptsize Planificar, Supervisar,\\ Diagnosticar (\textbf{12-B})};

			% Flujo nominal
			\draw[->, thick] (p1) -- (p2);
			\draw[->, thick] (p2) -- (p3);
			\draw[->, thick] (p3) -- (p4);
			\draw[->, thick] (p4) -- (p5);
			\draw[->, thick] (p5) -- (p6);
			\draw[->, thick] (p6) -- node[above]{\tiny \textbf{Sound}} (p7);
			\draw[->, thick] (p7) -- (p8);

			% Bucle de retroalimentación ante falla de verificación
			\draw[->, dashed, red!80, thick] (p6.north) to[out=135, in=-45] 
				node[above, sloped]{\tiny \alert{\textbf{Si falla: Regresar a 1 o 2}}} (p2.south);
		\end{tikzpicture}
	\end{center}
\end{frame}

% =======================================================
\section{Propiedades Formales de Comportamiento}
% =======================================================

\begin{frame}{}
	\centering
	\Huge{Propiedades Formales de Comportamiento: \\ ¿Qué le exigimos a un modelo correcto?}
\end{frame}

\begin{frame}{Alcanzabilidad y Estados Prohibidos}
	\begin{block}{Definición Formal de Alcanzabilidad (\emph{Reachability})}
		Un marcado $M$ es alcanzable desde el marcado inicial $M_0$ ($M \in R(M_0)$) si existe una secuencia finita de disparos $\sigma = t_1 t_2 \dots t_k$ tal que:
		\[ M_0 \xrightarrow{\;\sigma\;} M \]
		El conjunto $R(M_0)$ es el \textbf{Espacio de Estados} del sistema.
	\end{block}

	\begin{columns}
		\column{.55\textwidth}
		\begin{alertblock}{Verificación de Seguridad / Exclusión}
			Si definimos un conjunto de estados prohibidos o catastróficos $S_{\text{inseguro}}$ (ej: dos trenes en el mismo tramo de vía, o la liberación de producto sin control de calidad):
			\[ R(M_0) \cap S_{\text{inseguro}} = \emptyset \]
			Garantiza matemáticamente la imposibilidad física o lógica de acceder al error.
		\end{alertblock}

		\column{.40\textwidth}
		\begin{exampleblock}{Propiedad Dinámica}
			La alcanzabilidad depende intrínsecamente del marcado inicial $M_0$, no solo de la topología de la red.
		\end{exampleblock}
	\end{columns}
\end{frame}

\begin{frame}{Acotamiento (\emph{Boundedness}) y Seguridad (\emph{Safeness})}
	\begin{block}{Definición Formal de $k$-Acotamiento}
		Un lugar $p \in P$ es \textbf{$k$-acotado} si el número de fichas en $p$ nunca excede un entero $k$:
		\[ \forall M \in R(M_0): \quad M(p) \le k \]
		La red completa es $k$-acotada si todos sus lugares lo son.
	\end{block}

	\begin{columns}
		\column{.48\textwidth}
		\begin{alertblock}{Seguridad (\emph{Safeness} / 1-Acotada)}
			Si $k=1$, la red es **Safe (Segura)**:
			\[ \forall p \in P, \forall M \in R(M_0): \quad M(p) \le 1 \]
			\begin{itemize}
				\item Cada lugar representa una condición booleana pura (activa o inactiva).
				\item En workflows: garantiza que no hay tareas duplicadas accidentalmente en un mismo caso.
			\end{itemize}
		\end{alertblock}

		\column{.48\textwidth}
		\begin{exampleblock}{Significado en la Planta}
			\begin{itemize}
				\item Evita el desbordamiento de almacenes intermedios (\emph{buffers}).
				\item Previene la saturación de memoria en colas de software y bases de datos.
				\item Si una red no es acotada, el espacio de estados es **infinito**.
			\end{itemize}
		\end{exampleblock}
	\end{columns}
\end{frame}

\begin{frame}{Vivacidad (\emph{Liveness}) y Grados de Vivacidad}
	\begin{block}{La Jerarquía de Vivacidad (Niveles $L_0$ a $L_4$)}
		Una transición $t \in T$ tiene un grado de vivacidad según su capacidad de dispararse en el futuro:
		\begin{itemize}
			\item \textbf{$L_0$ - Muerta (\emph{Dead}):} $t$ nunca se puede disparar en ninguna secuencia alcanzable desde $M_0$. Representa \alert{código muerto o diseño defectuoso}.
			\item \textbf{$L_1$ - Cuasi-viva (\emph{Quasi-live}):} $t$ puede dispararse al menos una vez en alguna secuencia $\sigma$.
			\item \textbf{$L_2$ - Repetible $k$ veces:} Para cualquier entero $k$, existe una secuencia donde $t$ se dispara al menos $k$ veces.
			\item \textbf{$L_3$ - Infinitamente repetible:} Existe una secuencia de disparo infinita donde $t$ aparece infinitas veces.
			\item \textbf{$L_4$ - Completamente Viva (\emph{Live}):} Desde \textbf{cualquier} marcado $M \in R(M_0)$, $t$ sigue siendo potencialmente habilitable en algún paso futuro.
		\end{itemize}
	\end{block}

	\begin{alertblock}{Importancia en Automatización}
		Una red libre de bloqueos permanentes exige que sus actividades críticas sean vivas en un grado acorde al ciclo de vida del sistema.
	\end{alertblock}
\end{frame}

\begin{frame}{Bloqueos: \emph{Deadlock} vs. \emph{Livelock}}
	\begin{columns}
		\column{.48\textwidth}
		\begin{alertblock}{1. Bloqueo Total (\emph{Deadlock})}
			\begin{itemize}
				\item Se alcanza un marcado $M_d \in R(M_0)$ tal que:
				\[ \forall t \in T: \quad \neg (M_d [t\rangle) \]
				\item Ninguna transición de la red puede dispararse. El sistema queda congelado permanentemente.
				\item Típico en disputa circular de recursos: el robot A tiene la pieza y espera la máquina; la máquina tiene otra pieza y espera al robot A.
			\end{itemize}
		\end{alertblock}

		\column{.48\textwidth}
		\begin{block}{2. Bloqueo Dinámico (\emph{Livelock})}
			\begin{itemize}
				\item El sistema sigue disparando transiciones activamente.
				\item Sin embargo, queda confinado en un conjunto cerrado de estados cíclicos y \textbf{nunca progresa hacia el objetivo final} (ej: sumidero).
				\item En workflows: un bucle infinito de rechazos y reintentos sin condición de escape.
			\end{itemize}
		\end{block}
	\end{columns}
\end{frame}

% =======================================================
\section{Métodos de Verificación y Análisis}
% =======================================================

\begin{frame}{}
	\centering
	\Huge{Métodos de Verificación y Análisis: \\ Espacio de Estados, Álgebra y Topología}
\end{frame}

\begin{frame}{El Enfoque del Espacio de Estados: Árbol de Cobertura}
	\begin{block}{Árbol y Grafo de Alcanzabilidad}
		Se construye partiendo de $M_0$ disparando sucesivamente todas las transiciones habilitadas.
		\begin{itemize}
			\item Si la red es acotada, el grafo es **finito** y permite inspeccionar todas las propiedades directamente.
			\item \alert{Problema de la Explosión Combinatoria:} Una red pequeña con 15 lugares concurrentes puede generar millones de estados.
		\end{itemize}
	\end{block}

	\begin{columns}
		\column{.55\textwidth}
		\begin{exampleblock}{El Algoritmo de Karp-Miller y el Símbolo $\omega$}
			Para redes no acotadas, si en una rama se alcanza un marcado $M' \ge M$ ($M' \neq M$):
			\begin{itemize}
				\item Los lugares donde $M'(p) > M(p)$ crecen indefinidamente.
				\item Se reemplaza su valor por el símbolo $\omega$ (infinito).
				\item Produce un **Árbol de Cobertura Finito** que detecta formalmente si la red no está acotada.
			\end{itemize}
		\end{exampleblock}

		\column{.40\textwidth}
		\begin{center}
			\begin{tikzpicture}[node distance=1.1cm, >=Stealth, thick]
				\node[draw, rectangle, rounded corners] (n0) {\scriptsize $[1, 0]$};
				\node[draw, rectangle, rounded corners, below of=n0] (n1) {\scriptsize $[1, \omega]$};
				\draw[->] (n0) -- node[right]{\tiny $t_1$} (n1);
				\draw[->] (n1) edge [loop right] node[right]{\tiny $t_1$} (n1);
			\end{tikzpicture}
		\end{center}
	\end{columns}
\end{frame}

\begin{frame}{Análisis Algebraico: Invariantes de Lugar ($P$-Invariantes)}
	\begin{block}{Definición de $P$-Invariante}
		Un vector no nulo $y \ge 0$ de dimensión $|P|$ es un \textbf{$P$-invariante} si:
		\[ y^T \cdot C = 0 \]
		donde $C$ es la matriz de incidencia de la red.
	\end{block}

	\begin{alertblock}{Propiedad Fundamental de Conservación}
		Multiplicando la ecuación de estado por $y^T$:
		\[ y^T \cdot M = y^T \cdot M_0 + y^T \cdot C \cdot u = y^T \cdot M_0 = \text{Constante} \]
		\textbf{La suma ponderada de fichas en esos lugares permanece idéntica en todo momento.}
	\end{alertblock}

	\begin{exampleblock}{Interpretación Física en Plantas y Flujos}
		\begin{itemize}
			\item \textbf{Conservación de Recursos:} $M(\texttt{Robot\_Libre}) + M(\texttt{Robot\_Ocupado}) = 1$.
			\item \textbf{Acotamiento Estructural:} Si todo lugar pertenece a un $P$-invariante con pesos positivos, la red es \textbf{estructuralmente acotada} (independientemente de cuántos estados tenga).
		\end{itemize}
	\end{exampleblock}
\end{frame}

\begin{frame}{Análisis Algebraico: Invariantes de Transición ($T$-Invariantes)}
	\begin{block}{Definición de $T$-Invariante}
		Un vector no nulo $x \ge 0$ de dimensión $|T|$ es un \textbf{$T$-invariante} si:
		\[ C \cdot x = 0 \]
	\end{block}

	\begin{columns}
		\column{.50\textwidth}
		\begin{block}{Propiedad de Retorno al Estado}
			Si una secuencia realizable $\sigma$ tiene como vector de conteo de disparos a $x$:
			\[ M \xrightarrow{\;\sigma\;} M \]
			El sistema ejecuta la secuencia $\sigma$ y \textbf{regresa exactamente al mismo marcado inicial $M$}.
		\end{block}

		\column{.46\textwidth}
		\begin{exampleblock}{Aplicación en Procesos}
			\begin{itemize}
				\item Modela **ciclos completos de manufactura** o procesamiento de órdenes.
				\item Demuestra la **reversibilidad** y reproducibilidad de los ciclos operativos.
				\item Permite auditar si el flujo tiene caminos cíclicos no deseados.
			\end{itemize}
		\end{exampleblock}
	\end{columns}
\end{frame}

\begin{frame}{Estructuras Topológicas Críticas: Sifones y Trampas}
	\begin{columns}
		\column{.48\textwidth}
		\begin{alertblock}{Sifón (\emph{Siphon})}
			Un conjunto de lugares $S \subseteq P$ es un \textbf{sifón} si todas las transiciones que alimentan a $S$ son también de salida de $S$:
			\[ {}^\bullet S \subseteq S^\bullet \]
			\begin{itemize}
				\item \textbf{Propiedad Crítica:} Si en algún momento $S$ pierde todas sus fichas ($M(S)=0$), \alert{nunca más podrá recibir fichas}.
				\item Es la **génesis estructural del Deadlock**.
			\end{itemize}
		\end{alertblock}

		\column{.48\textwidth}
		\begin{exampleblock}{Trampa (\emph{Trap})}
			Un conjunto de lugares $Q \subseteq P$ es una \textbf{trampa} si todas las transiciones de salida de $Q$ son también de entrada a $Q$:
			\[ Q^\bullet \subseteq {}^\bullet Q \]
			\begin{itemize}
				\item \textbf{Propiedad Vital:} Si en algún momento $Q$ contiene al menos una ficha, \textbf{jamás quedará completamente vacío}.
				\item Protege al sistema contra la inanición.
			\end{itemize}
		\end{exampleblock}
	\end{columns}

	\vspace{0.2cm}
	\centering{\footnotesize \textbf{Teorema de Commoner (en redes Free-Choice):} La red es viva si y solo si todo sifón contiene una trampa marcada en $M_0$.}
\end{frame}

% =======================================================
\section{La Corrección de Flujos de Trabajo: Soundness}
% =======================================================

\begin{frame}{}
	\centering
	\Huge{La Corrección de Flujos de Trabajo: \\ Soundness y el Teorema del Cierre}
\end{frame}

\begin{frame}{Los Tres Criterios de Corrección (\emph{Soundness} - van der Aalst)}
	\begin{block}{Definición Formal de Soundness en una WF-Net}
		Una WF-Net $PN = (P, T, F)$ con entrada $i$ y salida $o$ es **Correcta (\emph{Sound})** si y solo si cumple:
		\begin{enumerate}
			\item \textbf{Opción de Completar (\emph{Option to complete}):}\\
			Para todo marcado alcanzable $M \in R([i])$, el marcado sumidero $[o]$ es alcanzable:
			\[ \forall M \in R([i]): \quad [o] \in R(M) \]
			\item \textbf{Finalización Adecuada (\emph{Proper completion}):}\\
			Si un marcado $M \in R([i])$ pone una ficha en el sumidero $o$, no puede haber fichas en ningún otro lugar:
			\[ \forall M \in R([i]): \quad (M \ge [o]) \implies (M = [o]) \]
			\item \textbf{Ausencia de Tareas Muertas (\emph{No dead tasks}):}\\
			Toda transición $t \in T$ puede dispararse en al menos una instancia de ejecución ($L_1$-viva):
			\[ \forall t \in T, \exists M \in R([i]): \quad M[t\rangle \]
		\end{enumerate}
	\end{block}
\end{frame}

\begin{frame}{Patologías Comunes: ¿Qué pasa cuando un WF no es \emph{Sound}?}
	\begin{columns}
		\column{.48\textwidth}
		\begin{alertblock}{1. Falta de Sincronización (Fichas Basura)}
			\begin{itemize}
				\item Se usa una bifurcación AND-split pero luego se recolecta con un OR-join en vez de AND-join.
				\item Una rama llega a $o$ y la otra rama queda ejecutándose sola.
				\item \textbf{Consecuencia:} Fichas residuales (\emph{garbage tokens}) contaminando la memoria y tareas disparadas dos veces por error.
			\end{itemize}
		\end{alertblock}

		\column{.48\textwidth}
		\begin{alertblock}{2. Sincronización Imposible (Deadlock)}
			\begin{itemize}
				\item Se usa una bifurcación excluyente (OR-split) pero luego se intenta consolidar con un AND-join.
				\item El AND-join espera fichas de ambas ramas, pero solo una se ejecutó.
				\item \textbf{Consecuencia:} Bloqueo absoluto. El caso de negocio queda colgado de por vida esperando una tarea que jamás ocurrirá.
			\end{itemize}
		\end{alertblock}
	\end{columns}
\end{frame}

\begin{frame}{El Teorema Fundamental del Cierre de van der Aalst}
	\begin{block}{La Red Cerrada $\overline{PN}$}
		Dada una WF-Net, conectamos el sumidero $o$ de regreso a la fuente $i$ mediante una transición virtual $t^*$:
	\end{block}

	\begin{center}
		\begin{tikzpicture}[node distance=2.2cm, >=Stealth, thick]
			\node[place, tokens=1] (i) [label={above:\scriptsize $i$ (Inicio)}] {};
			\node[draw, dashed, rectangle, minimum width=3.5cm, minimum height=1.2cm, right of=i, xshift=0.5cm] (wf) {Cuerpo del Workflow (WF-Net)};
			\node[place] (o) [right of=wf, xshift=0.5cm, label={above:\scriptsize $o$ (Fin)}] {};
			
			\draw[->] (i) -- (wf);
			\draw[->] (wf) -- (o);
			
			\node[transition] (tstar) [below of=wf, yshift=0.2cm, label={below:\scriptsize Transición de Cierre $t^*$}] {};
			\path (o) edge [->, bend left=45] (tstar);
			\path (tstar) edge [->, bend left=45] (i);
		\end{tikzpicture}
	\end{center}

	\begin{alertblock}{Teorema del Cierre (van der Aalst, 1997)}
		\[ \text{Una WF-Net es \textbf{Correcta (Sound)}} \iff \text{La red cerrada } \overline{PN} \text{ es \textbf{VIVA y ACOTADA}} \]
	\end{alertblock}
	\centering{\footnotesize \textbf{Trascendencia:} ¡No necesitamos una teoría nueva para flujos de trabajo! Todo el análisis algebraico y topológico de Redes de Petri verifica directamente cualquier workflow empresarial.}
\end{frame}

\begin{frame}{Validación en Tiempo de Carga en el Motor Fénix}
	\begin{block}{El Filtro de Despliegue (*Deploy Time*)}
		En la arquitectura de Fénix, la especificación en Red de Petri no es solo documentación gráfica, es un **contrato ejecutable**:
		\begin{enumerate}
			\item El diseñador modela el flujo de trabajo (en planta o administrativo).
			\item El compilador de Fénix cierra la red con $t^*$ y analiza:
				\begin{itemize}
					\item Invariantes de lugar $\implies$ acotamiento de recursos.
					\item Ausencia de sifones no cubiertos por trampas $\implies$ sin bloqueos.
					\item Vivacidad y acotamiento de $\overline{PN} \implies$ **Soundness garantizado**.
				\end{itemize}
			\item \textbf{Si el análisis pasa:} El flujo se despliega al motor para planificar y supervisar.
			\item \textbf{Si el análisis falla:} El motor rechaza el despliegue e indica exactamente cuál ruta conduce a bloqueo o a fichas huérfanas.
		\end{enumerate}
	\end{block}
\end{frame}

% =======================================================
\section{Conclusión}
% =======================================================

\begin{frame}{Conclusión de la Sesión 12-A}
	\begin{block}{Lo que aprendimos en esta sesión}
		\small
		\begin{itemize}
			\item \textbf{La Falacia de las Pruebas Empíricas:} La concurrencia exige demostración matemática formal previa; probar en producción es prohibitivo.
			\item \textbf{Propiedades Formales:} Alcanzabilidad ($R(M_0)$), Acotamiento y Seguridad ($M(p)\le 1$), y Vivacidad ($L_0$ a $L_4$).
			\item \textbf{Métodos de Análisis:} Árbol de cobertura con $\omega$ (Karp-Miller), $P$-invariantes (conservación física), $T$-invariantes (ciclos) y Sifones/Trampas (anatomía del bloqueo).
			\item \textbf{Soundness y Teorema del Cierre:} Opción de completar, finalización adecuada y sin tareas muertas $\iff \overline{PN}$ es viva y acotada.
			\item \textbf{Fénix Deploy-Time:} Garantía matemática de cero estados colgados antes de ejecutar.
		\end{itemize}
	\end{block}

	\begin{exampleblock}{Próxima Sesión (12-B): El Motor en Acción}
		\small
		**Planificación, Supervisión en Tiempo Real y Diagnóstico de Rendimiento**: ¿Por qué un estado demoró más de lo esperado? (Recursos vs. Calidad de Insumos vs. Variabilidad).
	\end{exampleblock}
\end{frame}

% Diapositiva final
\setbeamertemplate{headline}{}
\setbeamertemplate{background}
{\includegraphics[width=\paperwidth,height=\paperheight,keepaspectratio]{salida.pdf}}

\begin{frame}
	\placetextbox{0.6}{0.87}{\Huge\textbf{¡Muchas gracias!}}
\end{frame}

\setbeamertemplate{headline}{}
\setbeamertemplate{background}{}

\printbibliography
\end{document}
